% 勒文海姆–斯科伦定理（综述）
% license CCBYSA3
% type Wiki

本文根据 CC-BY-SA 协议转载翻译自维基百科\href{https://en.wikipedia.org/wiki/L\%C3\%B6wenheim\%E2\%80\%93Skolem_theorem}{相关文章}。

在数理逻辑中，洛文海姆–斯科伦定理（Löwenheim–Skolem theorem）是关于模型的存在性与基数的一个定理，以数学家莱奥波德·洛文海姆（Leopold Löwenheim）和托拉夫·斯科伦（Thoralf Skolem）的名字命名。  

其精确定义如下：若一个可数的一阶理论（countable first-order theory）存在一个无限模型，那么对于任意无限基数 \(\kappa\)，该理论都存在一个基数为 \(\kappa\) 的模型。这意味着，任何具有无限模型的一阶理论都不可能在同构意义下拥有唯一模型。  

由此可知，一阶理论无法控制其无限模型的基数。（向下的）洛文海姆–斯科伦定理是刻画一阶逻辑的两个关键性质之一，另一个是紧致性定理（compactness theorem）。这两者共同构成了林德斯特伦定理（Lindström’s theorem）的基础。而在更强的逻辑体系中（例如二阶逻辑），一般而言洛文海姆–斯科伦定理并不成立。[1]  
\subsection{定理}
\begin{figure}[ht]
\centering
\includegraphics[width=6cm]{./figures/9c693c47e7fbd008.png}
\caption{洛文海姆–斯科伦定理示意图（Illustration of the Löwenheim–Skolem theorem）} \label{fig_Skolem_1}
\end{figure}
在其最一般的形式下，洛文海姆–斯科伦定理（Löwenheim–Skolem Theorem）表述如下：

对于任意签名（signature）\(\sigma\)，任意无限的 \(\sigma\)-结构 \(M\)，以及任意满足 \(\kappa \geq |\sigma|\) 的无限基数 \(\kappa\)，存在一个 \(\sigma\)-结构 \(N\)，使得 \(|N| = \kappa\)，并且：
\begin{itemize}
\item 若 \(\kappa < |M|\)，则 \(N\) 是 \(M\) 的一个**初等子结构**（elementary substructure）；  
\item 若 \(\kappa \geq |M|\)，则 \(N\) 是 \(M\) 的一个**初等扩张**（elementary extension）。  
\end{itemize}
该定理通常被分为上述两种情况：  
第一部分（关于较小基数的初等子结构）称为向下洛文海姆–斯科伦定理（Downward Löwenheim–Skolem Theorem）[2]:160–162；第二部分（关于较大基数的初等扩张）称为向上洛文海姆–斯科伦定理（Upward Löwenheim–Skolem Theorem）[3]。  
\subsection{讨论}
以下将进一步阐述签名（signature）与结构（structure）的基本概念。
\subsubsection{概念}
\textbf{签名}

一个签名由以下部分组成：一组函数符号\(S_{\text{func}}\)，一组关系符号\(S_{\text{rel}}\)，以及一个函数\(\operatorname{ar} : S_{\text{func}} \cup S_{\text{rel}} \to \mathbb{N}_0\)表示各符号的元数（arity）。一个元数为 0 的函数符号称为常数符号。在一阶逻辑中，签名有时也被称为语言（language）。若其函数符号与关系符号的集合是可数的，则称该签名为可数签名。一般而言，签名的基数就是它所包含全部符号集合的基数。

一个一阶理论（first-order theory）由固定的签名与该签名下的一组句子（无自由变元的公式）组成[4]:40。理论常常通过：给出一组公理以生成该理论，或指定一个结构并令该理论由该结构满足的所有句子组成，  
来定义。

\textbf{结构 / 模型}

给定一个签名 \(\sigma\)，一个 \(\sigma\)-结构 \(M\) 是对 \(\sigma\) 中符号的具体解释。它包括一个基础集合（通常也记为 \(M\)），以及对函数符号和关系符号的解释。一个常数符号在 \(M\) 中的解释就是 \(M\) 的某个元素；一个 \(n\)-元函数符号 \(f\) 的解释是一个映射 \(f^M : M^n \to M\)；一个关系符号 \(R\) 的解释是 \(M\) 上的一个 \(n\)-元关系，即 \(M^n\) 的一个子集。

一个 \(\sigma\)-结构 \(M\) 的子结构（substructure）是这样获得的：  
取 \(M\) 的一个子集 \(N\)，使其对所有函数符号的解释封闭（因此包含所有常数符号的解释），并将关系符号的解释限制到 \(N\)。初等子结构（elementary substructure）是子结构中特殊的一种，它与原结构（其初等扩张）满足完全相同的一阶句子。
\subsubsection{推论}
在引言中给出的结论可直接由取 \( M \) 为该理论的任意无限模型而得。  
该定理的“向上”部分的证明还表明：若一个理论具有任意大的有限模型，则它必然也有一个无限模型；有时这一结果也被视为定理的一部分。[2]

一个理论若仅有一个模型（模同构而言），则称其为范畴的（categorical）。该术语最早由 Veblen（1904）提出。在随后的若干年间，数学家曾希望能够通过描述某种版本的集合论的一阶范畴理论，从而为整个数学建立一个坚实的逻辑基础。然而，洛文海姆–斯科伦定理给这一希望第一次致命打击：它表明，一个拥有无限模型的一阶理论不可能是范畴的。随后在 1931 年，哥德尔不完备定理（Gödel’s incompleteness theorem）彻底摧毁了这一希望。[2]

在 20 世纪早期，洛文海姆–斯科伦定理的许多推论在逻辑学家看来颇为违反直觉，因为当时人们尚未充分理解“一阶性质”与“非一阶性质”的区别。其中一个重要的推论是：存在不可数模型的“真算术”（true arithmetic），这些模型满足所有一阶归纳公理，但却含有非归纳子集。

设 \(\mathbb{N}\) 表示自然数集，\(\mathbb{R}\) 表示实数集。根据该定理可知：
理论 \((\mathbb{N}, +, \times, 0, 1)\)（即真的一阶算术理论）存在不可数模型；理论 \((\mathbb{R}, +, \times, 0, 1)\)（即实闭域理论）存在可数模型。当然，存在一些公理化系统能刻画 \((\mathbb{N}, +, \times, 0, 1)\) 与 \((\mathbb{R}, +, \times, 0, 1)\) 到同构为止，但洛文海姆–斯科伦定理说明：这些公理化系统不可能是一阶的。例如，在实数理论中，用于刻画 \(\mathbb{R}\) 为“完备有序域”的线性序的完备性性质，便是一个非一阶性质。[2]:161

另一个被认为尤其令人不安的推论是：存在一个可数的集合论模型，但该模型仍然满足“实数集是不可数的”这一命题。康托尔定理（Cantor’s theorem）指出某些集合是不可数的，而这一现象产生的悖论性局面被称为斯科伦悖论（Skolem’s paradox）。它揭示了一个深刻的事实：“可数性”这一概念并非绝对，而是相对模型而定的。[5]
\subsection{证明概述}
\subsubsection{向下部分}
对于每一个一阶 \(\sigma\)-公式\(\varphi(y, x_1, \ldots, x_n)\),选择公理（Axiom of Choice）蕴含存在一个函数  
\[
f_\varphi : M^n \to M~
\]
使得对所有 \(a_1, \ldots, a_n \in M\)，要么：
\[
M \models \varphi(f_\varphi(a_1, \ldots, a_n), a_1, \ldots, a_n)~
\]
要么：
\[
M \models \neg \exists y \, \varphi(y, a_1, \ldots, a_n)~
\]
再次应用选择公理，我们可以得到一个从所有一阶公式 \(\varphi\) 到这些函数 \(f_\varphi\) 的映射。

这一族函数 \(f_\varphi\) 诱导出一个作用于 \(M\) 的幂集上的**预闭包算子**（preclosure operator）：
\[
F(A) = \{\, f_\varphi(a_1, \ldots, a_n) \in M \mid 
\varphi \in \sigma, \; a_1, \ldots, a_n \in A \,\},
\quad A \subseteq M~
\]
对 \(F\) 进行可数次迭代得到一个闭包算子 \(F^\omega\)。取任意子集 \(A \subseteq M\)，满足 \(|A| = \kappa\)，定义\(N = F^\omega(A)\)。可以证明 \(|N| = \kappa\)。根据Tarski–Vaught 判别法，此时 \(N\) 是 \(M\) 的一个初等子结构。

此证明中使用的技巧最早可追溯到 斯科伦（Skolem），他在语言中引入了表示这些斯科伦函数\(f_\varphi\) 的新函数符号。我们也可以将 \(f_\varphi\) 定义为部分函数，仅在
\(M \models \exists y\, \varphi(y, a_1, \ldots, a_n)\)时有定义。  

关键之处在于，\(F\)是一个预闭包算子，它保证\(F(A)\)包含了所有以\(A\)中元素为参数并且在\(M\)中存在解的公式的解。同时满足：
\[
|F(A)| \le |A| + |\sigma| + \aleph_0~
\]
\subsubsection{向上部分}
首先，通过为 \(M\) 的每一个元素添加一个新的常数符号，扩展签名，得到扩展签名 \(\sigma'\)。在这个扩展签名下，\(M\) 的完整理论称为其初等图像（elementary diagram）。接着，向签名中再加入 \(\kappa\) 个新的常数符号，并向 \(M\) 的初等图像中添加句子\(c \neq c'\)（其中 \(c, c'\) 是任意两个不同的新常数符号）。根据紧致性定理（compactness theorem），该理论是一致的。由于其模型的基数至少为 \(\kappa\)，根据该定理的“向下部分”，我们可以得到一个恰好具有基数 \(\kappa\) 的模型 \(N\)。该模型 \(N\) 含有 \(M\) 的一个同构拷贝，且该拷贝是 \(N\) 的一个初等子结构。[6][7]:100–102

